\chapter{Wykład V}

\section{p-grupy Sylowa}
Niech	
\begin{itemize}
	\item $G$-grupa,
	\item $\left| G \right| = n$
	\item $k\in \N$
\end{itemize}
wtedy \[
\exists\left(H\le G\right)\left( \left| H \right| = k \overset{\text{wniosek z tw Lagrange'a}}{\Longrightarrow} k | n\right)  
.\] 
Zauważmy, że implikacji w drugą stronę tutaj nie mamy.

Kontrprzykład:

$6|12 = \left| A_4 \right| $ ale nie istnieje $H\le A_4$, takie że $\left| H \right| = 6$ ($\leftarrow$ zadanie)

Dążymy do częściowego odwrócenia tego wniosku, tzn dla szczególnych k- potęg liczb pierwszych.

Ustalmy p- liczba pierwsza.

\begin{definition}[p-grupa]
	G nazywamy p-grupą, gdy $\left| G \right| $ potęgą liczby pierwszej. Można założyć, że $p>0$.	
\end{definition} 
\begin{eg}
	\begin{itemize}
		\item $\mathcal{Q}_{8}$ jest 2-grupą, 
		\item $D_{2^{n}}$ jest 2-grupą,
		\item  $\Z_{p^{k}}$ jest $p$-grupą
		\item $\Z_{p}^{l}$ jest $p$-grupą 
			\[
			\underbrace{\Z_{p}\times \Z_{p}\times \ldots\times \Z_{p}}_{l}
			.\] 
	\end{itemize}
\end{eg} 
Będziemy analizować strukturę $p$-grup. W tym celu najpierw zbadamy własności działań $p$-grup na zbiorach.

\begin{definition}
	$G$ działa na $X$. Oznaczamy:
	\[
	\fix_{X}(G) := \{x\in X: \forall\left(g\in G\right)\left( g\cdot x = x \right)  \} 
	.\] 
	Jest to zbiór punktów stałych działania $G$ na $X$.
\end{definition} 
Gdy kontekst jest znany, to piszemy po prostu $\fix(G) $

\textbf{Przykłady} 
	
\begin{enumerate}[label=\arabic*.]
	\item $G$ działa na sobie przez automorfizmy wewnętrzne, wtedy 
		 \[
		\fix_{G}(G) = Z(G)
		.\] 
	\item $G\neq \{e\}$ działa na $G$ poprzesz lewe translacje, czyli $g\cdot x:= gx$, wtedy
		 \[
		\fix_{G}(G) = \O
		.\] 
	\item $GL_{n}(\R)$ działa na $\R^{n}$ w naturalny sposób, czyli $A\cdot x:= Ax$, wtedy
		\[
			\fix_{\R^{n}}(GL_{n}(\R)) = \{\vec{0}\} 
		.\] 
\end{enumerate}

\begin{context}

	Zakładamy teraz, że $G$ działa na $X$
\end{context}
Niech $x\in X$. Wtedy następujące warunki są równoważne:
\begin{enumerate}[label=\alph*)]
	\item $x\in \fix_{X}(G) $.
	\item $\left| G\cdot x \right| = 1$.
	\item $G\cdot x = \{x\} $
\end{enumerate}
\begin{proof}
	TODO
\end{proof} 
\begin{lemma}[O działaniach $p$-grup.]
	Niech $p$-grupa $G$ działa na $X$ i $\left| X \right| < \infty$. Wtedy $\left| \fix_{X}(G)  \right| \left| X \right|  \equiv \pmod{p}  $
\end{lemma} 
\begin{proof}
	Niech $x_1,x_2,\ldots,x_{n}\in X\setminus \fix_{X}(G)$ będą, takie że
	\[
		X = \fix_{X}(G) \uplus G\cdot x_1 \uplus \ldots \uplus G\cdot x_{n}\tag{$\uplus$ to suma rozłączna}
	.\] 
	Stąd 
	 \begin{align*}
		\left| X \right| &= \left| \fix_{X}(G)  \right| + \left| G\cdot x_1 \right| + \ldots + \left| G\cdot x_{n} \right|  \\
		&= \left| \fix_{X}(G)  \right| + \left[ G:Gx_{1} \right] + \ldots + \left[ G:x_{n} \right] 
	.\end{align*}
	
	Zauważmy, że $x_{i}\not\in \fix_{X}(G) \implies G\neq G_{x_{i}} \implies \left[ G:G_{x_{i}} \right] \neq 1$,
	ale $\left[ G:G_{x_{i}} \right] \mid \left| G \right| $, a $\left| G \right| $ jest potęgą liczby $p$.

	Stąd wynika $\left| X \right|  \equiv \left| \fix_{X}(G)  \right| \pmod{p} $
\end{proof} 

\begin{theorem}
	$G$-nietrywialna $p$- grupa $\implies p\mid \left| Z(G) \right| $w szczególności $Z(G)$ jest nietrywialne. 
\end{theorem} 
\begin{proof}
	Stosujemy dowód lematu dla $G$ działajacej na sobie przez automorifzmy wewnętrzne. Wtedy
	\[
	\left| Z(G) \right| = \left| \fix_{X}(G)  \right| \equiv \left| G \right|  \equiv 0 \pmod{p} \implies p\mid \left| Z(G) \right|  
	.\] 
\end{proof} 

\textbf{Przykłady}
\begin{enumerate}[label=\arabic*.]
	\item $-\mathcal{I}\in Z(\mathcal{Q}_{8}).$ 
	\item $O_{\pi}\in Z(D_{2^{n}}).$
\end{enumerate}
\begin{lemma}
	Niech $\varphi : G \to H$ będzie homomorfizmem, $N\le H$. Wtedy
	\begin{enumerate}[label=\arabic*.]
		\item Jeśli $G$ jest skończona i $\varphi $ jest ,,na'', to 
			\[
			\left| \varphi ^{-1} \left[ N \right]  \right| = \left| \ker(\varphi )  \right| \cdot \left| N \right| .\]
		\item $N\unlhd H \implies \varphi ^{-1}\left[ N \right] \unlhd G$
	\end{enumerate}
\end{lemma} 
\begin{proof}
	TODO
\end{proof} 
\begin{theorem}
	Niech $\left| p^{m} \right|.$ Wtedy istnieje ciąg podgrup
	\[
	G_0\le G_1\le \ldots\le G_{m} = G
	,\] 
	taki że
	\[
		\forall\left(i\right)\left( \left| G_{i} \right| = p^{i}\text{ oraz }G_{i}\unlhd G \right)  
	.\] 
\end{theorem} 
\begin{proof}
	Indukcja względem m.
\end{proof} 

\begin{corollary}
Twierdzenie Lagrange'a można odwrócić dla $p$-grup. 
\end{corollary}
\begin{definition}
	Niech $H\le G$ oraz $G$ jest skończona.
	\begin{enumerate}[label=\arabic*.]
		\item $H$ jest $p$-podgrupą gdy $H$ jest $p$-grupą.
		\item $H$ jest $p$-podgrupą Sylowa, gdy $H$ jest $p$-podgrupą i $p\nmid \left[ G:H \right] = \frac{|G|}{|H|} $.
		\item Przez $\mathfrak{X}_{p}$ oznaczamy zbiór wszystkich $p$-podgrup Sylowa grupy $G$.
		\item $s:= \left| \mathfrak{X}_{p} \right| $
	\end{enumerate}
\end{definition} 
\begin{remark}
	$\left| G \right| = p^{n\cdot k}$ dla pewnych $n,k\in \N$, takich że $p\nmid k$. Piszemy wtedy $p^{n}\parallel \left| G \right| $. Czytamy to, że $p^{n}$ dokładnie dzieli $\left| G \right| $. Wtedy $H\le G$ jest $p $-podgrupą Sylowa $\iff \left| H \right| = p^{n}$.
\end{remark} 

\textbf{Przykłady}

\begin{enumerate}[label=\arabic*)]
	\item $G = \Z_{12}$, $\left| G \right| = 12 = 2^{2}\cdot 3$
		\begin{itemize}
			\item 2-podgrupy Sylowa mają moc 4. Jest jedyna taka: $\{0,3,6,9\} $, czyli $s_{2} = 1$ 
			\item 3-podgrupy Sylowa mają moc 3. Jest jedyna taka:  $\{0,4,8\} $, czyli $s_3 = 1$.

				(w grupie cyklicznej zawsze jest jedyna?)
		\end{itemize}
	\item $G=S_3$, $\left| G \right| = 6 = 2\cdot 3$ 
		\begin{itemize}
			\item 2-podgrupy Sylowa są mocy 2: $\langle \left( 1,2 \right)  \rangle , \langle \left( 1,3 \right)  \rangle , \langle \left( 2,3 \right)  \rangle $.

				Stąd wynika że $s_{2}= 3$
			\item 3-podgrupy Sylowa: $\langle \left( 1,2,3 \right)  \rangle = A_3 \implies s_{3} = 1$
		\end{itemize}
	\item $G = A_4$, $\left| G \right| = 12 = 2^2\cdot 3$
		\begin{itemize}
			\item 2-podgrupy Sylowa są rzędu 4: $\{\id, \left( 1,2 \right)\left( 3,4 \right), \left( 1,3 \right)\left( 2,4 \right), \left( 1,4 \right)\left( 2,3 \right) \} \implies s_2=1$ 
			\item 3-podgrupy Sylowa są rzędu 4:
				$\langle \left( 1,2,3 \right)  \rangle , \langle \left( 1,2,4 \right)  \rangle , \langle \left( 1,3,4 \right)  \rangle , \langle \left( 2,3,4 \right)  \rangle \implies s_3= 4$
		\end{itemize}
\end{enumerate}
\begin{theorem}[Sylowa]
	Załóżmy, że $\left| G \right| = p^{n}\cdot k, n\ge 1\text{ oraz }p\nmid k$. Wtedy
	\begin{enumerate}[label=\arabic*.]
		\item Zawsze istnieje  $p$-podgrupa Sylowa H grupy G ($\implies k = \left[ G:H \right] $ ) 
		\item Dla każdej $p$-podgrupy grupy $G$ i dla każdego $H\in \mathfrak{X}_{p}$ istnieje $g\in G$, $K\subseteq gHg^{-1}$:
			\[
				\forall\left(K \colon \text{$p$-podgr. } G\right) \forall\left(H \in \mathfrak{X}_{p}\right) \exists\left(g \in G\right) \left( K \subseteq gHg^{-1} \right)
			.\] 
		\item $\forall\left(H_1,H_2\in \mathfrak{X}_{p}\right)\exists\left(g\in G\right)\left( gH_1g^{-1} = H_2\right)   $ 
		\item $s_{p} \mid \left| G \right| $ oraz $s_{p} \equiv 1 \pmod{p}$ ($\implies s_{p}\mid k$) oraz $s_{p} = \left[ G: N_{G}(H) \right] $, dla $H \in \mathfrak{X}_{p}$
	\end{enumerate}
\end{theorem} 
W powyższym twierdzeniu są zawarte 1,2 i 3 twierdzenia Sylowa.
\begin{proof}
	Będziemy używać wielu różnych działań grup.
	\begin{enumerate}[label=\arabic*.]
		\item Niech $X:= \{A\subseteq G: \left| A \right| = p^{n}\} $. Mamy tu zbiór podzbiorów $G$ mocy $p^{n}$.
			$\left| X \right| = \begin{pmatrix} p^{n}\cdot k \\ p^{n} \end{pmatrix} $ TODO
	\end{enumerate}
\end{proof} 

\begin{theorem}[Twierdzenie Cauchy'ego]
	Załóżmy, że $p\mid \left| G \right| $. Wtedy $\exists\left( g\in G \right)\left( \text{rząd}\left( g \right) = p\right)  $
\end{theorem} 
\begin{proof}
	TODO
\end{proof} 
\begin{remark}
	Twierdzenia Cauchy'ego nie można udowodnić do wszystkich potęg liczby $p$, na przykład $4=2^2 \mid \Z_2\times \Z_2$ nie ma elementu rzędu 4.
\end{remark}
\begin{theorem}
	Niech $H$ będzie $p $- podgrupą Sylowa grupy G. Wtedy $H\unlhd G \iff s_{p}=1$ (czyli H jest jedyną $p$-podgrupą Sylowa)
\end{theorem} 
\begin{proof}
	TODO
\end{proof} 
\begin{theorem}
	Dla $G$ będącą grupą skończoną $p$-podgrupa $H$ grupy $G$ jest $p$-podgrupą Sylowa wtedy i tylko wtedy gdy $H$ jest maksymalną $p$-podgrupa
	(w sensie inkluzji)
\end{theorem} 
\begin{proof}
	TODO
\end{proof} 
\begin{eg}
	$G=A_4$, $\left| A_4 \right| = 2^2\cdot 3$, czyli 2-podgrupy Sylowa są rzędu 4, 3-podgrupy Sylowa są rzędu 4.
	W $A_4$ jest 8 elementów rzędu 3(cykle rzędu 3) i dają łącznie 4 podgrupy rzędu 3, czyli $s_{3} = 4$

	Pozostają 4 elementy: $\id,\left( 1,2 \right)\left( 3,4 \right) , \left( 1,3 \right) \left( 2,4 \right) , \left( 1,4 \right)\left( 2,3 \right) $ 
	czyli muszą stanowić jedynie 2-podgrupy Sylowa $H$. Z powyższego wniosku $\implies H\unlhd G$
\end{eg} 
\begin{remark}	
	W $A_4$ mamy $s_2=1$, $s_3=4$ i mamy dzielnik normalny rzędu 4.
	\begin{question}	
		Czy można to dostać z samego twiedzenia Sylowa i tego, że $\left| G \right| =12$?
	\end{question}
	$s_2 \equiv 1 \pmod{2} $ oraz $s_2\mid 3\implies s_2=1 \lor s_2=3$,

	$s_3 \equiv 1 \pmod{3} $ oraz $s_3\mid 4\implies s_3=1 \lor s_3=4$

	\begin{enumerate}[label=(\roman*)]
		\item dla $G=\mathbb{Z}_{12}$ : $s_2=1$ i $s_3=1$,
		\item dla $G=S_3\times \mathbb{Z}_{2}$: $s_2=3$ i $s_3=1$ 
		\item (ZADANIE) Nie moze sie zdarzyć $s_2=3$ i $s_3=4$
	\end{enumerate}
\end{remark} 
\textbf{Przykłady} 

\begin{enumerate}[label=\alph*)]
	\item $\left| G \right| = 40 = 2^3\cdot 5$. $s_5\mid 8$i $ s_5 \equiv 1 \pmod{5} \implies s_5=1$. Stąd $G$ ma dzielnik normalny rzędu 5.
	\item $\left| G \right| = 30 = 2\cdot 3\cdot 5$. Pokażemy, że istnieje $N\unlhd G$ takie że $\left| N \right| = 3 $lub $\left| N \right| = 5$ 
		\begin{proof}
			TODO
		\end{proof} 
	\item $\left| G \right| = 24 = 2^3\cdot 3$. Pokażemy, że istnieje $N\unlhd G$, taka że $N\neq \{e\} $ oraz $N\neq G$
		\begin{proof}
			TODO na 2 sposoby
		\end{proof} 
\end{enumerate}
\begin{remark}
	Niech $H\le G$. Wtedy $\bigcap_{g\in G} gHg^{-1}$ jest najmniejszym dzielnikiem normalnym.
\end{remark} 
\begin{proof}
	TODO (jest jako zadanie)
\end{proof} 
